\subsection{Chunking e fusione gerarchica}

Un ultimo problema, trasversale a tutti gli stadi precedenti, è quello del \emph{chunking}: un documento specie se lungo, come un regolamento o una dispensa, non può essere indicizzato come unico blocco, sia per limiti di
lunghezza degli input dei modelli di embedding sia perché un frammento troppo esteso diluisce il segnale specifico
cercato dalla domanda in mezzo a informazione non pertinente. La scelta della dimensione dei frammenti è però essa
stessa un compromesso: frammenti troppo corti rischiano di perdere il contesto necessario a renderli comprensibili
in isolamento, frammenti troppo lunghi tornano a soffrire del problema opposto.

La strategia più semplice consiste nel suddividere il testo in blocchi di dimensione fissa, eventualmente con una
sovrapposizione tra blocchi consecutivi per non spezzare un'informazione a cavallo del confine. Questo approccio tuttavia ignora la
struttura logica del documento: un paragrafo coerente può essere tagliato a metà solo perché il conteggio dei token ha raggiunto la soglia impostata. Un'alternativa più sensibile al contenuto è il \emph{chunking strutturale}, che rispetta i confini naturali del documento (sezioni, paragrafi, righe di una tabella) invece di tagliare a una lunghezza fissa, cedendo una parte dell'uniformità dei frammenti risultanti in cambio di un testo internamente più coerente. Nessuna delle due strategie, da sola, risolve però il problema opposto: un frammento strutturalmente coerente ma troppo corto (un singolo titolo di sezione, una riga isolata di una tabella) resta comunque privo del contesto che lo renderebbe comprensibile se recuperato in isolamento.

Una soluzione a questo problema è la \emph{fusione gerarchica} dei frammenti. In fase di
indicizzazione ogni documento viene diviso in frammenti piccoli, che sono quelli effettivamente
indicizzati e confrontati con la domanda, ma si conserva anche il legame di ciascun frammento con
la sezione del documento a cui appartiene. Al momento del recupero, se tra i frammenti selezionati
ne compaiono abbastanza appartenenti alla stessa sezione, questi vengono sostituiti dalla sezione
intera: il recupero conserva così la precisione dei frammenti piccoli, mentre il modello riceve un
contesto ampio e coerente.
